% =============================================================================
% Appendices -- NOT counted against the 40--60 page limit.  Use them.
%
% Under the narrative-first structure, the appendices carry the derivations,
% entry-level formulae, per-case numbers and software detail that the body
% points to.  Self-containment rule: every number quoted in the body must be
% derivable from the body or from here -- never from the repository.
% =============================================================================

\section{Quantum Gate Reference}
% Resolves the \todo in §2.1.2.

\section{Chebyshev and Ultraspherical Identities}
\label{app:cheb}
% Entries of Btilde and Dtilde; the banded differentiation recurrences.
% Referenced from §4.3.

This appendix records the entries of the banded factors $\tilde{B}$ and
$\tilde{D}$ introduced in \Cref{sec:descriptor-construction}, the exact factorisation
$G = \tilde{B}^{-1}\tilde{D}$, and the fully populated $N = 4$ instance of the
first-order descriptor operator.

\subsection{The differentiation recurrence}

For $f = \sum_{\ell\ge 0} a_\ell T_\ell$ with derivative $f' = \sum_{\ell\ge 0} p_\ell T_\ell$ on
$[-1,1]$, the Chebyshev coefficients satisfy
\begin{equation}
    c_{\ell-1}\,p_{\ell-1} - p_{\ell+1} = 2\ell\,a_\ell , \qquad \ell \ge 1 ,
    \label{eq:app-cheb-recurrence}
\end{equation}
with $c_0 = 2$ and $c_\ell = 1$ for $\ell \ge 1$; the $\ell = 1$ row therefore carries the
factor $c_0 = 2$ on $p_0$. Truncating to $N = 2^n$ coefficients and dropping the
$p_N$ term, \Cref{eq:app-cheb-recurrence} is the matrix identity
$\tilde{B}_{\mathrm{std}}\,p = \tilde{D}_{\mathrm{std}}\,a$ with
$\tilde{D}_{\mathrm{std}}$ a single superdiagonal, $(\tilde{D}_{\mathrm{std}})_{\ell,\ell+1} = 2(\ell+1)$,
and $\tilde{B}_{\mathrm{std}}$ two diagonals,
$(\tilde{B}_{\mathrm{std}})_{\ell,\ell} = c_\ell$, $(\tilde{B}_{\mathrm{std}})_{\ell,\ell+2} = -1$.

\subsection{The normalised basis}

The construction of \Cref{sec:descriptor-construction} works in the normalised basis
$\varphi_0 = T_0 / 2^{n/2}$, $\varphi_\ell = T_\ell / 2^{(n-1)/2}$ ($\ell \ge 1$), orthonormal
under the discrete $N$-point Chebyshev inner product $\langle T_k, T_\ell \rangle_N =
\tfrac{N}{2}\,c_\ell\,\delta_{k\ell}$ \cite{trefethen2000spectral}: since $N = 2^n$, this gives
$\|\varphi_0\|_N = \|\varphi_\ell\|_N = 1$ exactly, so the coefficient vector of a unit-norm
function is a unit vector. Writing
$S = \operatorname{diag}(\sqrt{c_\ell})$, the normalised factors are
$\tilde{B} = S\,\tilde{B}_{\mathrm{std}}\,S^{-1}$ and
$\tilde{D} = S\,\tilde{D}_{\mathrm{std}}\,S^{-1}$, that is
\begin{equation}
    \tilde{D}_{\ell,\ell+1} = 2(\ell+1)\sqrt{c_\ell / c_{\ell+1}} , \qquad
    \tilde{B}_{\ell,\ell} = c_\ell , \qquad
    \tilde{B}_{\ell,\ell+2} = -\sqrt{c_\ell / c_{\ell+2}} ,
    \label{eq:app-normalised-entries}
\end{equation}
so $\tilde{D}$ has the single non-zero diagonal $(2\sqrt{2},\, 4,\, 6,\, 8,\, \dots)$
and $\tilde{B}$ has main diagonal $(2,\, 1,\, 1,\, \dots)$ and second superdiagonal
$(-\sqrt{2},\, -1,\, -1,\, \dots)$. Both hold $\mathcal{O}(N)$ non-zero entries at
any $N$. The dense Chebyshev derivative matrix is their exact quotient,
\begin{equation}
    G = \tilde{B}^{-1}\tilde{D} ,
\end{equation}
an algebraic identity rather than a numerical approximation: $\tilde{B}$ and $\tilde{D}$
are read directly off the same recurrence that defines $G$
(\Cref{eq:app-cheb-recurrence}), so the equality holds exactly at every $N$,
independent of whichever values a test suite happens to check.
\Cref{app:verification} reports this identity's
residual against the reference implementation, alongside the rest of the
composition-algebra checks.

\subsection{The Ultraspherical Connection}
\label{app:ultraspherical}

The banded factorisation above is the first-derivative case of the
ultraspherical spectral method of Olver and Townsend \cite{olver2013}. Stating
the correspondence precisely is what fixes what this thesis adds to it.

Olver and Townsend represent a function's derivative not in the Chebyshev
basis $T_k$ but in the ultraspherical (Gegenbauer) basis $C_k^{(\lambda)}$,
the family of polynomials orthogonal with respect to the weight
$(1-x^2)^{\lambda - 1/2}$; at $\lambda = 1$ these are the Chebyshev
polynomials of the second kind, $C_k^{(1)} = U_k$. Differentiation of a
Chebyshev series obeys
\begin{equation}
    \frac{\mathrm{d}T_k}{\mathrm{d}x} = k\,C^{(1)}_{k-1} , \qquad k \ge 1
    \label{eq:app-ot-diff}
\end{equation}
\cite[Eq.~2.5]{olver2013}, so for $u = \sum_k u_k T_k$ the coefficients of
$u'$ in the $C^{(1)}$ basis are read off directly: $(\mathcal{D}_0 u)_k =
(k+1)u_{k+1}$, a single superdiagonal, sparse at every truncation. Because
$u'$'s coefficients now live in a different basis to $u$'s own, a second,
purely algebraic operator is needed to convert a function's own Chebyshev
coefficients into the $C^{(1)}$ basis without differentiating it:
$\mathcal{S}_0$, with $(\mathcal{S}_0)_{k,k} = 1$ ($k=0$), $\tfrac12$ ($k \ge
1$), and $(\mathcal{S}_0)_{k,k+2} = -\tfrac12$ \cite[Eq.~2.8]{olver2013}.
Comparing entries, $\tilde{D}_\mathrm{std} = 2\mathcal{D}_0$ and
$\tilde{B}_\mathrm{std} = 2\mathcal{S}_0$ exactly: the recurrence
$\tilde{B}_\mathrm{std}\,p = \tilde{D}_\mathrm{std}\,a$ of
\Cref{eq:app-cheb-recurrence} is literally $\mathcal{S}_0\,p = \mathcal{D}_0\,a$,
doubled by the convention that folds the weight of $T_0$ into $c_0 = 2$
rather than into the polynomial itself. This is the discrete shadow the
framing of \Cref{sec:descriptor-construction} claims: the descriptor
recurrence does not compute a new identity, it asserts an existing one --
that $p$'s own Chebyshev coefficients, converted into the $C^{(1)}$ basis by
$\mathcal{S}_0$, must equal the derivative computed directly by
$\mathcal{D}_0$ -- as one row of a linear system, rather than solving it for
$p$ via $\mathcal{S}_0^{-1}$.

The generalisation to order $k \ge 1$ in \cite{olver2013} raises the basis
one level per derivative: a $\lambda$-th derivative needs coefficients in
$C^{(\lambda)}$, obtained from the analogous pair
$\mathrm{d}C_k^{(\lambda)}/\mathrm{d}x = 2\lambda\,C_{k-1}^{(\lambda+1)}$
($k\ge1$, giving $\mathcal{D}_\lambda$, still a single superdiagonal) and a
conversion operator $\mathcal{S}_\lambda : C^{(\lambda)} \to C^{(\lambda+1)}$,
still two diagonals \cite[Eqs.~3.2--3.3]{olver2013}. To assemble one equation
of order $k$, every term is raised to the common top basis $C^{(k)}$ by a
chain $\mathcal{S}_{k-1}\cdots\mathcal{S}_{\lambda+1}$ applied after each
term's own $\mathcal{M}_\lambda[a^\lambda]\mathcal{D}_\lambda$, giving one
almost-banded operator $\mathrm{L} = \mathcal{M}_k[a^k]\mathcal{D}_k +
\sum_{\lambda=0}^{k-1}\mathcal{S}_{k-1}\cdots\mathcal{S}_\lambda
\mathcal{M}_\lambda[a^\lambda]\mathcal{D}_\lambda$ that is solved once for
the coefficients of $u^{(k)}$ \cite[\S3.2]{olver2013} -- so a $k$-th order
equation needs $k$ distinct conversion operators $\mathcal{S}_0, \dots,
\mathcal{S}_{k-1}$ chained together and the lower-order unknowns eliminated.

The descriptor construction of \Cref{sec:descriptor-construction} takes a
different route to the same sparsity. Rather than raising the basis and
eliminating, it keeps every derivative's coefficients in the \emph{same}
normalised Chebyshev basis, as its own block $c^{(j)}$ of $w$, and reuses the
single $\lambda = 0$ pair $(\tilde{B}, \tilde{D})$ between every consecutive
pair of blocks (\Cref{eq:descriptor-recurrence-rows}). No elimination is
performed, and no operator beyond $\tilde{B}$, $\tilde{D}$ and the
multiplication operators of \Cref{sec:app-variable-coeff} is required at any
order $k$ -- the price is the larger state $w$ that carries every
intermediate derivative explicitly instead of eliminating it
(\Cref{sec:app-variable-coeff}). What this thesis adds to \cite{olver2013}
is not a new sparsification of differentiation -- that is their result -- but
applying the \emph{same} $\lambda = 0$ identity block-wise to build an
unreduced, differential-algebraic residual for the physics-informed
Hamiltonian of \cite{wu2025pihm}, whose block structure is exactly what
\Cref{sec:descriptor-atoms} block-encodes.

\subsection{The \texorpdfstring{$N = 4$}{N = 4} instance}

At $n = 2$ ($N = 4$), \Cref{eq:app-normalised-entries} gives
\begin{equation}
    \tilde{B} =
    \begin{pmatrix}
        2 & 0 & -\sqrt{2} & 0 \\
        0 & 1 & 0 & -1 \\
        0 & 0 & 1 & 0 \\
        0 & 0 & 0 & 1
    \end{pmatrix},
    \qquad
    \tilde{D} =
    \begin{pmatrix}
        0 & 2\sqrt{2} & 0 & 0 \\
        0 & 0 & 4 & 0 \\
        0 & 0 & 0 & 6 \\
        0 & 0 & 0 & 0
    \end{pmatrix},
    \qquad
    G =
    \begin{pmatrix}
        0 & \sqrt{2} & 0 & 3\sqrt{2} \\
        0 & 0 & 4 & 0 \\
        0 & 0 & 0 & 6 \\
        0 & 0 & 0 & 0
    \end{pmatrix}.
\end{equation}
The top-right entry $3\sqrt{2}$ of $G$ is the fill-in: coefficient $0$ of the
derivative depends on coefficient $3$ of the input. Row $\ell$ of $G$ is
non-zero at columns $\ell+1, \ell+3, \dots$, so this is the extreme case --
the lowest-order coefficient can depend on up to half the others, while the
highest-order one (row $3$ above) depends on none. The factorisation removes
the fill-in entirely: no row of $\tilde{B}$ or $\tilde{D}$ touches more than
two coefficients.

For the first-order equation $f' + 2f = 0$ of \Cref{sec:descriptor-construction},
$w = (c^{(0)};\, c^{(1)})$ and
$A = \left(\begin{smallmatrix} -\tilde{D} & \tilde{B} \\ 2 I_4 & I_4 \end{smallmatrix}\right)$
is the $8 \times 8$ matrix
\begin{equation}
    A =
    \left[
    \begin{array}{cccc|cccc}
        0 & -2\sqrt{2} & 0 & 0 & 2 & 0 & -\sqrt{2} & 0 \\
        0 & 0 & -4 & 0 & 0 & 1 & 0 & -1 \\
        0 & 0 & 0 & -6 & 0 & 0 & 1 & 0 \\
        0 & 0 & 0 & 0 & 0 & 0 & 0 & 1 \\ \hline
        2 & 0 & 0 & 0 & 1 & 0 & 0 & 0 \\
        0 & 2 & 0 & 0 & 0 & 1 & 0 & 0 \\
        0 & 0 & 2 & 0 & 0 & 0 & 1 & 0 \\
        0 & 0 & 0 & 2 & 0 & 0 & 0 & 1
    \end{array}
    \right] .
    \label{eq:app-descriptor-8x8}
\end{equation}
The top four rows are the recurrence block $-\tilde{D}\,c^{(0)} + \tilde{B}\,c^{(1)} = 0$; the
bottom four are the differential row $2c^{(0)} + c^{(1)} = 0$. Every non-zero entry lies on a
diagonal or a near-diagonal, in place of the dense $G + 2I$ of the eliminated
form.

\subsection{Higher order, variable coefficients, and the padded block count}
\label{sec:app-variable-coeff}

For an order-$k$ equation the state is $w = (c^{(0)};\, \dots;\, c^{(k)})$ and
$A$ stacks $k$ recurrence block-rows
$-\tilde{D}\,c^{(j)} + \tilde{B}\,c^{(j+1)} = 0$ above the single differential row
$\sum_j M_{a_j}\,c^{(j)} = 0$: $k+1$ block-rows, each $N$ coefficients wide,
giving the classical $(k+1)N \times (k+1)N$ square matrix of
\Cref{def:descriptor-residual-operator}.

A variable coefficient $a_j(x)$ enters through $M_{a_j} = \sum_q c_q\,M_{x^q}$,
$0 \le q \le N$, where $c_q$ are the \emph{monomial} coefficients of $a_j$: supplied
directly as an array $[c_0, c_1, \dots]$, or, when $a_j$ is given as a function,
obtained by a least-squares fit to degree $N-1$ in the Chebyshev basis followed
by an exact Chebyshev-to-monomial change of basis. The Chebyshev fit is what
pins $q$ to the same resolution $N$ already fixed for the equation; its degree,
not the subsequent change of basis, is the only place $a_j$'s representation is
truncated. $M_{x^q}$ is the
resolution-$n$, square multiplication-by-$x^q$ operator built from
$x\,T_\ell = \tfrac{1}{2}(T_{\ell+1} + T_{|\ell-1|})$ with the $T_N$ term dropped. $M_{x^q}$
is banded with bandwidth exactly $q$, so it folds into the same shift-and-diagonal
form as $\tilde{B}$ and $\tilde{D}$. The dropped $T_N$ term is the only
approximation in the linear construction: for a smooth product resolved below
degree $N - q$ the resulting error is confined to the top $\sim q$ Chebyshev modes
and measures $\sim 10^{-13}$.

\paragraph{Padding to a power of two.} The classical $(k+1)N \times (k+1)N$
matrix above is exact and complete as a linear system; it is not, however,
the object any later stage of the pipeline actually builds. A gate-level
realisation must index which block $c^{(j)}$ a given operator acts on with a
qubit register, and a register of $m$ qubits addresses $2^m$ blocks, not
$k+1$ of them for an arbitrary $k$. $A$ is therefore \emph{embedded} in a
$2^{\lceil \log_2(k+1) \rceil}N$-dimensional space: the $k+1$ block-rows and
block-columns above are kept exactly as derived, and
$2^{\lceil \log_2(k+1) \rceil} - (k+1)$ further diagonal blocks are appended,
each an identity $I_N$ pinning its own dummy field to zero,
$c^{(j)}_{\text{dummy}} = 0$. This padding is exact: the physical solution already satisfies the dummy rows
trivially, since it has nothing in those blocks, so no error is introduced
beyond the $M_{x^q}$ truncation above. The qubit cost, however, is not
uniform in $k$:

\begin{center}
\begin{tabular}{ccccc}
\toprule
Order $k$ & Blocks $k+1$ & Field qubits $\lceil\log_2(k+1)\rceil$ & Padded blocks $2^{\lceil\log_2(k+1)\rceil}$ & Dummy blocks \\
\midrule
1 & 2 & 1 & 2 & 0 \\
2 & 3 & 2 & 4 & 1 \\
3 & 4 & 2 & 4 & 0 \\
4 & 5 & 3 & 8 & 3 \\
7 & 8 & 3 & 8 & 0 \\
\bottomrule
\end{tabular}
\end{center}

so the state $w$ the gate-level construction actually carries has
$2^{\lceil \log_2(k+1) \rceil}N$ components, not $(k+1)N$: equal to the
unpadded count exactly when $k+1$ is itself a power of two (orders $1, 3, 7,
\dots$), and up to nearly double it otherwise -- order $2$, the most common
case after order $1$, already pays for a field register it only three-quarters
uses. This is the qubit cost that must always accompany the gate-count
comparison of \Cref{sec:descriptor-cost}.

\section{The Descriptor Construction in Full}
% Atom-by-atom gate-level constructions and the composition algebra.
% Referenced from §4.4 and §4.5.

\section{Proof of the Cost Proposition}
% Referenced from §4.6.

\section{The Nonlinear Route in Full}
% Carleman lift, the nodal fold, doubled-space residual algebra.
% Referenced from §5.

\section{Verification Protocol and Residuals}
\label{app:verification}
% Per-atom residuals, composition-algebra tests, the tolerances quoted in §4.10.
% TODO: report the banded-factorisation residual ||B̃^{-1}D̃ - G|| here, backed
% by a pytest parametrize that actually spans the n range quoted -- see
% test_descriptor_encoding.py::test_banded_factorisation_is_exact.

\section{Catalogue of Differential Equation Cases}
% Every case solved, with parameters and status.  Referenced from §6.4--§6.6.

\section{Readout Protocols and Their Costs}
% Candidate protocols and sample costs quoted from the literature.
% Referenced from §4.9.  Nothing here was implemented in this work.

\section{Software}
% DESolverLib / desolver_hpc architecture, Setonix configuration, run parameters.

\section{Research Proposal}
% Required by the School guidelines; note deviations from the proposal.
